% legends.tex (v6): figure legends (each under 350 words) and Table 1 title.
\section*{Figure legends}

\textbf{Figure 1. Study regions, observations and time coverage.} (a) Permafrost probability (ESA CCI 1997--2021 mean) with the three regions outlined and the number of 1 km cells with observations. (b--d) Observed cells (dots) and cells observed in four or more years (stars) in Alaska, the Lena Delta and north-western Canada. (e) Cell-years per observation year and region; the text gives the share of the last observation month. (f) Years covered by the labels, the yearly ERA5-Land forcing, the satellite inputs and the compared products; the shaded band is the 2010--2024 analysis window.

\textbf{Figure 2. Thaw depth, degree-days and the evaluation design.} (a) Observed thaw depth within one season at two Lena Delta sites against the square root of the degree-days accumulated to the measurement day; lines are the Stefan fits with coefficient $E$. (b) Observed active-layer thickness against the square root of observation-day degree-days of that year for cells observed in several years; lines connect years of the same cell. (c) The validation splits on the Lena Delta cells: random by cell, spatial block (0.5$^\circ$) and year hold-out; colours mark the scoring fold.

\textbf{Figure 3. Time matching, validation splits and year-wise behaviour.} (a) Change in spatial-block RMSE when year-specific forcing replaces climatological forcing, by method and region; bars are 95\% block-bootstrap intervals. (b) RMSE of each method under random, spatial-block, year, block-by-year and region hold-out splits. (c) Correlation between predicted and observed interannual anomalies under year hold-out. (d) Observed against predicted year-to-year change at cells observed in several years (physics plus residual model). (e) Year-wise bias of the physics-plus-residual model under the spatial-block split for years with at least 50 cell-years. P0 source-coefficient Stefan, P1 recalibrated Stefan, Ku Kudryavtsev, D0 direct learning, F1 physics-input learning, R1 physics plus residual learning.

\textbf{Figure 4. Conditions of the gain and comparison with products.} (a) Change in spatial-block RMSE of the physics-plus-residual model relative to recalibrated Stefan along the covariate ladder. (b) The same change within regional terciles of thawing degree-days. (c) The same change with different learners substituted for CatBoost. (d) RMSE on the cell-years common to all compared products; filled markers use all common cells, open markers exclude cells within 25 km of the training sites of Wei et al. and Liu et al.; values above 45 cm are printed. (e) Decile means of predicted against observed thickness with interquartile bands.

\textbf{Figure 5. Yearly series at the longest-observed sites.} (a--f) Two sites per region: observations (open circles), recalibrated Stefan and physics-plus-residual predictions for every year 2010--2024 with end-of-season forcing (lines), the residual model at the observation day (squares) and the yearly values of Wei et al. and Liu et al. (g) Standard deviation of predicted against observed thickness across years at cells with four or more observed years. (h) Regional mean anomaly of the predicted thickness over all cells and of the observations for years with at least 50 cell-years.

\textbf{Figure 6. Year-labelled 1 km maps.} (a) Mean thickness 2011--2024 from the physics-plus-residual model over Alaska with observed cells (dots); faded cells lie outside the area of applicability. (b, c) Difference of the 2014 and 2022 maps from the 2011--2024 mean, each forced with the ERA5-Land degree-days of that year. (d) 1 km detail in the North Slope and Interior windows with observed cells circled. (e) Area of applicability. (f) Lena Delta thickness for 2016. (g) Samoylov--Kurungnakh window with observed cells.

% ---------------------------------------------------------------------------------------------------------------
% v7 figure set (outputs/figures/paper/v7, scripts/4_visualization/v7_*.py; 2026-10-15). Compiled only when
% \vsevenfigs is defined in main.tex (see the commented line next to \graphicspath). F4_results_core replaces
% Figure 3a-b and Figure 4a-c above and adds the model-structure panel (e). Legends for the other v7 figures are
% added by their figure tasks. Numbers are macros from numbers.tex (point estimate [95% interval]).
\ifdefined\vsevenfigs
\textbf{Figure 4. Year-specific forcing, validation splits and the conditions of the residual gain.} (a) Change in spatial-block RMSE when the forcing of the observation year replaces climatological forcing, by method and region; points are the change and bars the 95\% block-bootstrap interval. The region legend applies to (a) and (c--f); regions are offset within each category. (b) RMSE of four methods under random, spatial-block, year, block-by-year and region hold-out splits, one panel per region. Lines in (b) are coloured by method, not by region: source-coefficient Stefan (grey), recalibrated Stefan (blue), physics plus residual ML (orange) and direct ML (magenta, dashed). Under region hold-out there are no in-region observations to recalibrate, so recalibrated Stefan reverts to the source coefficients and the two Stefan lines meet. The open marker at the upper edge of the Lena Delta panel marks a value beyond the axis range (direct ML, region hold-out, \nrmseLeregDz{} cm). (c) Change in spatial-block RMSE of the physics-plus-residual model relative to recalibrated Stefan for six covariate sets: climate only; with terrain; with terrain and soil (the main set); the main set with MODIS vegetation and snow-off date; the main set with high-resolution soil and land-cover classes; and all covariates including lagged forcing. (d) The same change within regional terciles of thawing degree-days. (e) The same change for four model structures: direct ML (no physics), physics-input ML (Stefan prediction as a feature), physics pseudo-labels (training set augmented with unobserved cells labelled by the training-fold Stefan prediction) and physics plus residual ML (Stefan anchor plus a learned residual). With all labels, pseudo-labels change RMSE by \ndrefLeblkDo{} cm in the Lena Delta, \ndrefCablkDo{} cm in Canada and \ndrefAkblkDo{} cm in Alaska, reproducing the physical model without adding to it; the residual structure gives \ndrefLeblkRo{} and \ndrefCablkRo{} cm in the Lena Delta and Canada. (f) The same change with CatBoost and four alternative learners: random forest (RF), multilayer perceptron (MLP), tabular ensemble MLP (TabM) and feature-tokenizer transformer (FT-T); filled markers physics plus residual ML, open markers direct ML. In (c--f) bars are 95\% block-bootstrap intervals and negative values favour the learned model.
\fi

\textbf{Figure 7. Year-to-year maps, regional series and comparison with products.} (a) Alaska 1 km active-layer thickness anomaly of each year 2011--2024 from the physics-plus-residual model: the map of that year minus the 2011--2024 mean of the same cell, each year forced with the ERA5-Land degree-days of that year. The colour range is symmetric, $\pm$15 cm, the same for every year; values beyond the range take the end colours. Hatching marks cells outside the area of applicability (the same static mask in every year); grey land is not mapped (permafrost probability below 50\%, water, or missing inputs). Circles mark 0.5$^\circ$ blocks with observed cells in that year, sized by the number of cells. (b) Yearly anomaly of the mean over all mapped cells (line) for Alaska and the Lena Delta, with the mean anomaly of the observed cells (circles) for years with at least 50 cell-years; the observed values are centred on the predicted series of the same years. The Lena Delta 2020 value is labelled with its number of cell-years. (c) RMSE on the cell-years shared by all compared products and lying more than 25 km from the product training sites, for the physics-plus-residual model and the best product of each region. (d) Share of mapped cells inside the area of applicability and half-width of the 90\% prediction interval for the two mapped regions.

\section*{Table}

\textbf{Table 1. Observations, cells and years by region.} Number of cell-years, 1 km cells, 0.5$^\circ$ blocks, observation years, cells with four or more observed years, and the year with the largest share of cell-years. Sources and selection rules are given in Methods.
